\section{Related Work}
% \subsection{}
% Recent advances in using large language models (LLMs) for automated CTF solving have led to the development of agentic systems capable of reasoning, tool use, and decision making. Early efforts relied on single-agent pipelines, while newer frameworks adopt modular designs with task decomposition and role-based agents~\cite{guo2024largelanguagemodelbased,song2024audit,saha2025malgen,liu2024multi,bianou2024pentestmitre,dorri2018multi, abramovich2024enigma}. A recent work, D-CIPHER~\cite{udeshi2025d}, presents a more comprehensive multi-agent system that enhances coordination, improves task execution, and boosts overall performance through task delegation, feedback loops, and specialized agent roles.

\subsection{LLM Agents for Offensive Security}
Recent advancement of LLM agentic capabilities have led to the development of a variety of LLM agents for cybersecurity automation, relying on single-agent pipelinesto modular designs with task decomposition and role-based agents~\cite{guo2024largelanguagemodelbased,song2024audit,saha2025malgen,liu2024multi,bianou2024pentestmitre,dorri2018multi, abramovich2024enigma}. A recent work, \texttt{D-CIPHER}~\cite{udeshi2025d}, presents a more comprehensive multi-agent system that enhances coordination, improves task execution, and boosts overall performance through task delegation, feedback loops, and specialized agent roles.
Despite these advances, systematic evaluation of hyperparameter sensitivity remains an understudied area in the context of LLM agents. Prior works often evaluate with default parameters such as temperature and top-p, or tune them in isolation without exploring their joint effects on reasoning quality, execution success, or inter-agent communication. Some studies have acknowledged the impact of decoding behavior~\cite{turtayev2024hacking,yang2023intercode,yao2022react}, but none have conducted targeted investigations on how these hyperparameters influence performance in multi-agent CTF-solving scenarios. 
%This gap limits our understanding of optimization strategies for agentic systems and motivates the need for rigorous ablation and tuning studies in structured cybersecurity tasks.

\subsection{CTF Benchmarks} 
CTF competitions have served as challenging tasks to evaluate the cybersecurity capabilities of automated LLM agents. The NYU CTF Bench~\cite{shao2024nyu} and CyBench~\cite{zhang2024cybenchframeworkevaluatingcybersecurity} provide multi-category challenge sets and step-by-step annotations, respectively, enabling structured evaluation. However, these datasets are often broad in scope or require significant computational resources, making them less suitable for rapid experimentation. Moreover, few are tailored for evaluating the nuanced behaviors of agentic systems under varying configurations. To address this, we introduce \texttt{CTFTiny}, a compact  50-challenge benchmark that enables rapid experiments, reproducible baselines, and  exploration of agent design and parameter tuning strategies.
% \subsection{Multi-Agent Architectures in LLM Systems}

\subsection{LLM as a Judge}
Traditional CTF evaluation often reduces agent performance to a pass/fail metric. To understand nuanced progress, recent work has explored using LLMs as automated judges~\cite{gu2024survey,li2024llms,cao2025multi}. These LLM judges assess reasoning quality, partial progress, and the correctness of intermediate steps, providing a deeper understanding of agent capabilities beyond simple flag capture. In this work, we use an LLM-judge that offers granular scoring of agent behavior across multiple aspects of task execution.
% \nanda{Need to refine}